% =====================================================================
%  Capítulo 13 — Guardia y distancia
% =====================================================================
\chapter{Guardia y distancia}

Antes de ejecutar un golpe es importante estar equilibrada: columna estirada, con el peso de la parte superior dentro de la base que forman las piernas. Postura relajada, peso repartido entre los dos pies separados aproximadamente la anchura de los hombros, uno adelantado apuntando al frente y el otro atrasado a 45 grados hacia fuera. Con el peso sobre las puntas de los pies (asentarse en los talones desequilibra), flexionamos rodillas y ajustamos el tórax hasta movernos sin tambalearnos.

Es importante practicar esto hasta adoptar la posición instintivamente ante un peligro. Aunque la posición de los brazos en la calle no ha de ser agresiva, sí ha de estar atenta a bloquear: una \enquote{guardia de paz} con las palmas abiertas a la altura del pecho, en gesto de calma (\enquote{tranquilo, no quiero problemas}), protege sin provocar y deja las manos listas. Para entrenar dispondremos una guardia de mano abierta similar a la de los boxeadores, con el brazo no dominante adelante: el brazo adelantado protege desde la boca hasta el abdomen y el atrasado está listo para lanzar el ataque.

De pie frente a un compañero, medimos la distancia a la que debemos estar para llegar de manera efectiva con nuestras diferentes armas corporales. Cada persona tiene una distancia diferente según su anatomía, y debemos acostumbrarnos a no lanzar golpes que no van a impactar:

\begin{itemize}
  \item Distancia larga: patadas.
  \item Distancia media: manos.
  \item Distancia corta: codos y rodillas.
  \item Cuerpo a cuerpo.
\end{itemize}

En defensa personal entrenaremos sobre todo la distancia media, corta y el cuerpo a cuerpo, dejando las patadas altas a las artes marciales (requieren práctica habitual y comprometen el equilibrio). Hasta acostumbrarnos a golpear, atención a no estirar los miembros por completo: deben quedar imperceptiblemente retraídos para no dañar las articulaciones con cada golpe.

\section{Golpes de distancia media: la mano}

Si no tenemos entrenamiento en deportes de contacto, golpearemos con la mano abierta, utilizando la palma y su base (dedos semiextendidos y apretados): es más segura para nuestros huesos que el puño. Partiendo de la guardia, proyectamos la mano hacia la cabeza del oponente en un movimiento circular, teniendo en cuenta:

\begin{itemize}
  \item Acompañar cada golpe con la exhalación del aire.
  \item Corregir la distancia: ni tan corta que apenas toque, ni tan larga que el brazo no tenga recorrido.
  \item Comenzar sin potencia ni velocidad hasta familiarizarse con la técnica; después incrementar siempre dentro de la comodidad y sin dolor.
\end{itemize}

Cada golpe debe realizarse en equilibrio, pies bien asentados, acompañando cada mano con su giro de cadera. Por último, se practicarán estos golpes desde una actitud no agresiva (posición de descanso o \enquote{guardia de paz}) y de forma sorpresiva, sin movimientos previos que delaten la intención.

\section{Golpes de distancia corta: el codo}

El codo posee potencia y dureza naturales aun en personas que nunca han practicado artes de combate. Partimos de la guardia a distancia corta. Como con la mano, son primordiales el equilibrio y el giro de cadera: siempre será más fuerte el codazo de la pierna retrasada, al contar con mayor trayectoria. Proyectamos el codo hacia delante paralelo al suelo (el puño termina horizontal delante del pecho) impactando con la zona vecina a la punta.

Tras el golpe, la cintura queda colocada para sacar el otro codo girando la cadera en sentido opuesto, encadenando codazos. Con entrenamiento percibiremos esta mecánica y su progresiva potencia. Se sugiere su uso a muy corta distancia y para sorprender con un fuerte golpe a la mandíbula que provocará la confusión suficiente para seguir encadenando. Dominado el golpe horizontal, conviene entrenarlo en otras direcciones: ascendente a la mandíbula, descendente, y hacia atrás contra agarres por la espalda. Los codos también son muy recomendables para interponerse en la trayectoria de golpes y patadas.

\section{Golpes de distancia corta: la rodilla y golpes bajos}

La rodilla es un arma fuerte y dura por naturaleza: impactaremos en el ángulo que nos sea cómodo (ascendente al frente, o circular de fuera a dentro en corta distancia). Tendrá más efecto el rodillazo de la pierna atrasada por la mayor acción de la cadera. Objetivos: genitales y piernas y, según nuestra flexibilidad y la altura del oponente (por ejemplo, si le hemos bajado la cabeza tirando del pelo), estómago, costillas y cara.

Otras técnicas de nivel bajo se inician elevando poco la rodilla, sin anunciar el ataque, y estirando rápidamente la pierna:

\begin{itemize}
  \item Puntapié con la punta del zapato a tibia y rótula.
  \item Patada con la cara interna o externa del pie a tibia y rodilla, en trayectoria recta, como \enquote{cortando} la pierna.
  \item Pisotón al empeine y dedos del pie.
  \item Talonazo hacia atrás, para aflojar agarres desde la espalda.
\end{itemize}

Este entrenamiento, unido al uso de tacones, puede provocar daños de gran consideración por su capacidad de penetración, sobre todo \enquote{clavando} el tacón hacia delante o en pisotón.

\section{Sueltas básicas de agarres}

\begin{itemize}
  \item \textbf{Agarre de muñeca:} gira tu brazo y tira con fuerza hacia el hueco entre su pulgar y sus dedos (el punto débil de todo agarre), acompañando con todo el cuerpo, no solo el brazo.
  \item \textbf{Agarre de pelo:} fija sus manos contra tu cabeza con tus dos manos (quita el dolor y el control), gira hacia él y contraataca a ojos o garganta.
  \item \textbf{Abrazo de oso por detrás:} baja el centro de gravedad, talonazo a empeine/tibia, cabezazo hacia atrás si te aprisiona los brazos, y palanca a un dedo en cuanto liberes una mano.
  \item \textbf{Estrangulamiento frontal:} protege la vía aérea metiendo la barbilla, agarra y desgarra uno de sus meñiques mientras atacas a los ojos; gira el cuerpo para deshacer la presión.
\end{itemize}

Estas descripciones son solo recordatorios para quien las haya practicado: las sueltas requieren práctica supervisada con compañero.
